\documentclass[12pt]{article}
\usepackage[utf8]{inputenc}
\usepackage[T1]{fontenc}
\usepackage{hyperref}
\title{North Slavey Language Resources for K-12 Curriculum Development}
\author{Compiled by Claude Sonnet 5.5 (claude-sonnet-5-5) for Andr\'as Kornai}
\date{}
\begin{document}
\maketitle

\section{Introduction}
North Slavey (ISO 639-3 scs) is spoken in the Sahtu region of the Northwest Territories, Canada. Community names for it include Dene K\textschwa d\textschwa{} and Sahtu Dene. The 2000 Sahtu regional language plan found that speakers did not consider "North Slavey" or "Sahtu Dene" fully appropriate, so it uses "Dene" generally. It uses Sahtu\'ot'ine, K'ashog\'ot'ine and Sh\'uta\k{o}t'ine for the local dialect groups. The searches found catalogue entries, dictionaries, an app and the regional curriculum framework. They did not find a complete graded school curriculum in the language.

\section{Primary Educational Resources}
\begin{itemize}
\item Dene Kede Education: A Dene Perspective, Grades K--6 (GNWT, 1993), listed in the Sahtu Renewable Resources Board (SRRB) Dene Ked\textschwa{} catalogue. It is a cultural and language curriculum framework developed with elders and teachers from the five Dene regions. \url{https://www.srrb.nt.ca/people-and-places/dene-ked-to-dene-ts-l-catalogue-of-resources/item/2865-dene-kede-education-a-dene-perspective-grades-k-6}
\item Dene God\textschwa{} -- North Slavey Bibliography, published by the Sahtu Divisional Education Council (SDEC). It catalogues North Slavey language and literacy materials from 1979 to 2007. \url{https://srrb.nt.ca/people-and-places/dene-ked-to-dene-ts-l-catalogue-of-resources/item/3013-dene-god-north-slavey-bibliography}
\item North Slavey Alphabet Chart, in the SRRB catalogue. \url{https://www.srrb.nt.ca/people-and-places/dene-ked-to-dene-ts-l-catalogue-of-resources/item/2854-north-slavey-alphabet-chart}
\item Shutaot'ine Intro, a mobile learning app with lessons, games and quizzes. It was launched in 2016 by NWT Education, Culture and Employment, the Yamozha Kue Society and SDEC. \url{https://srrb.nt.ca/people-and-places/dene-ked-to-dene-ts-l-catalogue-of-resources/item/2970-shutaot-ine-intro}
\end{itemize}
Search results also reported that the Chief Jimmy Soldat School in Deline piloted instruction in the Dene language from kindergarten to grade 2 starting in 1979. They also reported a 1978 pre-primer, \emph{Sahtu Got'ine Gokede: A Slavey Language Pre-Primer in the Speech of Fort Franklin}. I did not find a link to the pre-primer itself.

\section{Linguistic and Reference Materials}
\begin{itemize}
\item Sh\'uta\k{o}t'ine Gok\textschwa de (Tulita Dictionary). It is a topical English-to-Slavey dictionary published by SDEC and open access via Google Drive. It is a revision of a 1986 Deline Teaching and Learning Centre work. It has no audio. \url{https://knowledgebase.arts.ubc.ca/?p=2952}
\item Record of the same dictionary on the UBC wiki. \url{https://wiki.ubc.ca/Documentation:RelLex/Sh%C3%BAta%C7%ABt%27ine_Gok%C9%99de_(Tulita_Dictionary)}
\item K'ash\'ogot'\k{i}ne Rad\k{i}l\k{i}h K\'o\'e X\textschwa d\textschwa\'{} (Fort Good Hope Dictionary). It is a topical wordlist of the Hare dialect, also published by SDEC, with no audio. \url{https://knowledgebase.arts.ubc.ca/8543-2}
\item Languagegeek, North Slavey. It is a general language overview page. \url{https://www.languagegeek.com/dene/kashogotine/5north_slavey.html}
\item Dene Speech Atlas, North Slavey Groups (University of Rochester). \url{https://sas.rochester.edu/lin/DeneSpeechAtlas/dene-language-groups/north-slavey/}
\item Wikipedia, Slavey language. \url{https://en.wikipedia.org/wiki/Slavey_language}
\end{itemize}

\section{Religious and Community Resources}
\begin{itemize}
\item The four Gospels translated into the Slave language (1883), a historical translation by Bompas and Reeve, digitized by Canadiana. It uses an older orthography and an older, probably South Slavey, dialect. \url{https://www.canadiana.ca/view/oocihm.14248}
\item The Gospel of St. Mark translated into the Slave language (1886), also digitized by Canadiana. \url{https://www.canadiana.ca/view/oocihm.15556}
\end{itemize}
No modern North Slavey scripture translation turned up in the searches.

\section{Digital and Audio Resources}
The Shutaot'ine Intro app, listed above, is the only digital learning tool found. The two dictionaries have no audio. The search results did not surface any audio or video archive.

\section{Curriculum Development Context}
The GNWT's Dene Kede curriculum has been in place since 1993, with grade 7--9 documents published from 2002 onward. Sahtu developers named in the catalogue include Jane Modeste, Albertine Baton and Therese Pellitier. SDEC is the regional body producing language materials. The 2000 Sahtu Dene language planning report is catalogued at the SRRB. \url{https://srrb.nt.ca/people-and-places/dene-ked-to-dene-ts-l-catalogue-of-resources/itemlist/tag/Language%20?start=60}

The NWT Department of Education, Culture and Employment page found in the searches may be relevant for contacts and programme status. \url{https://www.ece.gov.nt.ca/en/node/1824}

A curriculum would need graded readers, audio, and a fuller dictionary and grammar, because those were not found. Developers should work with SDEC and the community language authorities.
\end{document}